\chapter{Optimización en dimensión finita}

\noindent
Para el desarrollo de este tema, conviene tener en cuenta a la hora de resolver un problema de optimización en dimensión finita: % // TODO: Quizás desarrollar más este punto
\begin{itemize}
    \item Análisis de derivadas primeras y segundas: formas cuadráticas y polinomio de Taylor.
    \item Complicaciones en la frontera.
    \item Optimización con restricciones de igualdad\footnote{En el caso de las restricciones de tipo desigualdad tenemos los multiplicadores KKT (Karush-Kuhn-Tucker).}, método de los multiplicadores de Lagrange.
    \item Nociones de convexidad permiten pasar de nociones locales a globales.
    \item Derivadas direccionales.
    \item Programación lineal\footnote{Hay problemas similares, conocidos como programación cuadrática y programación convexa.}.
    \item Teorema de Weierstrass para la existencia de óptimos, previa a buscarlos.
\end{itemize}

\begin{teo}[de Weierstrass]
    Sea $K\subset \mathbb{R}^d$ un conjunto compacto y $f:K\to \mathbb{R}$ una aplicación continua. Entonces $f$ alcanza su mínimo.
    \begin{proof}
        Tomando $\alpha = \inf\limits_{x\in K}f(x)$, sabemos que existe una sucesión $\{x_n\}$ de puntos de $K$ de forma que $\{f(x_n)\}\to \alpha$. Por compacidad ha de existir $x\in K$ de forma que $\{x_n\}\to x$, y por continuidad de $f$ tenemos que $\{f(x_n)\}\to f(x)$, de donde $f(x) = \alpha$.
    \end{proof}
\end{teo}

\begin{definicion}[Semicontinua inferiormente]
    Sea $X$ un espacio topológico 1AN, decimos que una aplicación $f:X\to \mathbb{R}$ es semicontinua inferiormente si para toda sucesión convergente $\{x_n\}$ de puntos de $X$ se tiene que:
    \begin{equation*}
        \liminf \{f(x_n)\} \geq f\left(\lim \{x_n\}\right)
    \end{equation*}
\end{definicion}

\begin{teo}[Método directo del cálculo de variaciones]
    Sea $K$ un espacio topológico 1AN compacto y $F:K\to \mathbb{R}$ una aplicación semicontinua inferiormente. Entonces $F$ alcanza su mínimo.
    \begin{proof}
        Tomando $\alpha = \inf\limits_{x\in K}F(x)$, sabemos que existe una sucesión $\{x_n\}$ de puntos de $K$ de forma que $\{F(x_n)\}\to \alpha$. Por compacidad ha de existir $x\in K$ de forma que $\{x_n\}\to x$. Tenemos así que:
        \begin{equation*}
            \alpha = \lim \{F(x_n)\} = \liminf \{F(x_n)\} \geq F(x)
        \end{equation*}
        Y como $\alpha = \inf\limits_{x\in K}F(x)$ tendrá que ser $F(x) = \alpha$.
    \end{proof}
\end{teo}

\chapter{Optimización en dimensión infinita}
\section{Derivada de Gateaux}
\noindent
En dimensión infinita tenemos también el concepto de derivadas direccionales.\\

\noindent
Sea $X$ un espacio normado, $D\subset X$ y sea $F:D\to \mathbb{R}$. Fijamos $y\in D$, diremos que $\varphi \in X\setminus \{0\}$ es \textbf{dirección admisible para $y$ en $D$} si existe $\varepsilon_0>0$ tal que $y+\varepsilon\varphi \in D$ para\footnote{Podríamos definirlo en $\left]0,\varepsilon_0\right[$, pero por motivmos pedagógicos nos restringiremos a este caso, por simplicidad a la hora de dar los enunciados.} $\varepsilon \in \left]-\varepsilon_0,\varepsilon_0\right[$. Por tanto, está definida la correspondencia:
\begin{equation*}
    \varepsilon\longmapsto F(y+\varepsilon\varphi), \qquad \varepsilon \in \left]-\varepsilon_0,\varepsilon_0\right[
\end{equation*}
Supongamos que es derivable, en cuyo caso diremos que $F$ es derivable en el sentido de Gateaux en $y$ a lo largo de $\varphi$. La derivada correspondiente se define como:
\begin{equation*}
    \delta_\varphi F(y) = \lim_{\varepsilon\to0} \frac{F(y+\varepsilon\varphi)-F(y)}{\varepsilon} = \left[\frac{d}{d\varepsilon}F(y+\varepsilon\varphi)\right]_{\big|\varepsilon=0}
\end{equation*}
Diremos que $F$ es derivable en el sentido de Gateaux en un punto $y\in D$ si $F$ es derivable en el sentido de Gateaux en $y$ para toda dirección admisible $\varphi$.

\noindent
Diremos que $y_\ast\in D$ es un \textbf{punto crítico} (o estacionario) de $F$ si\footnote{También incluimos aquí los puntos en los que no hay ninguna dirección admisible.}:
\begin{equation*}
    \delta_\varphi F(y_\ast) = 0 \qquad \forall \varphi \text{\ dirección admisible}
\end{equation*}

\begin{ejemplo}
    Consideremos $D = \{C^1[-1,1] : y(-1) = 1 = y(1)\}\subset C^1[-1,1]$. Para $y(x) = x^2 \in D$, ¿son direcciones admisibles las siguientes?
    \begin{align*}
        \varphi_1(x) &= x^2 \\
        \varphi_2(x) &= 1-x^2
    \end{align*}
    Para que sean direcciones admisibles necesitamos comprobar si
    \begin{equation*}
        y_{\varepsilon}(x) = y(x) + \varepsilon\varphi_i(x)\in D \qquad \text{para $\varepsilon \in \left]-\varepsilon_0,\varepsilon_0\right[$}, \quad i \in \{1,2\}
    \end{equation*}
    \begin{enumerate}
        \item Para $y_\varepsilon(x) = x^2+\varepsilon x^2 = x^2(1+\varepsilon)$, tenemos que $y_{\varepsilon}(-1) = 1+\varepsilon\neq 1$, por lo que no es una dirección admisible para ningún $\varepsilon$.
        \item Para $y_{\varepsilon}(x) = x^2+\varepsilon (1-x^2)$ tenemos que:
            \begin{align*}
                y_{\varepsilon}(-1) &= 1 + \varepsilon(1-1) = 1 \\
                y_\varepsilon(1) &= 1 + \varepsilon(1-1) = 1
            \end{align*}
            Por lo que sí que es una dirección admisible.
    \end{enumerate}
    Cualquier $\varphi \in C^1[-1,1]$ con $\varphi(-1) = 0 = \varphi(1)$ será una dirección admisible.
\end{ejemplo}

\begin{notacion}
    Notaremos usualmante:
    \begin{equation*}
        C^1_0([a,b]) = \{y\in C^1([a,b]) : y(a) = 0 = y(b)\}
    \end{equation*}
\end{notacion}

\begin{prop}
    Sea $X$ un espacio normado, sea $D\subset X$ y $F:D\to \mathbb{R}$. Si $y_0\in D$ es extremo relativo de $F$ y $F$ es derivable de Gateaux en $y_0$, entonces $y_0$ es un punto crítico.
    \begin{proof}
        Supongamos que en $y_0$ se alcanza un mínimo local (en caso de máximo es análogo), es decir, existe un entorno de $y_0$ en el cual $F(y_0)\leq F(y)$ para cualquier $y$ en dicho entorno. Sin pérdida de generalidad, podemos suponer que es de la forma $D\cap B(y_0,\delta)$, para $\delta\in \mathbb{R}^+$.\\

        \noindent
        Dada $\varphi \in X$ dirección admisible para $y_0$, definimos $y_{\varepsilon} = y_0 + \varepsilon\varphi$, con lo que existe $\varepsilon_0>0$ tal que $y_\varepsilon \in D\quad \forall \varepsilon \in \left]-\varepsilon_0,\varepsilon_0\right[$. Observemos que:
        \begin{equation*}
            \|y_0 - y_{\varepsilon}\| = |\varepsilon| \|\varphi\| \quad\Longrightarrow\quad y_{\varepsilon}\in B(y_0,\delta) \quad \forall \varepsilon \in \left]\frac{-\delta}{\|\varphi\|}, \frac{\delta}{\|\varphi\|}\right[
        \end{equation*}
        Sea $\overline{\varepsilon} = \min\left\{\varepsilon_0,\frac{\delta}{\varphi}\right\}$, la correspondencia $\varepsilon\longmapsto F(y_{\varepsilon})$ está definida en $\left]-\overline{\varepsilon},\overline{\varepsilon}\right[$, es derivable y tiene un mínimo local en $\varepsilon=0$. En dicho caso, como $F$ es derivable de Gateaux en $y_0$ respecto a $\varphi$, obtenemos que $\delta_\varphi F(y_0) = 0$. Como la dirección admisible era arbitraria, se cumple para cualquiera.
    \end{proof}
\end{prop}

\section{Estudio de funcionales integrales}
\noindent
Tomaremos $X = C^1(I)$, con $I=[x_0,x_1]$ para $x_0< x_1$. Usaremos la norma:
\begin{equation*}
    \|y\|_{C^1} = \|y\|_\infty + \|y'\|_\infty
\end{equation*}
Consideraremos funcionales dados por: 
\begin{equation*}
    \cc{F}(y) = \int_{x_0}^{x_1} F(x,y(x),y'(x))~dx 
\end{equation*}
Típicamente consideraremos\footnote{O quizás en otro dominio de definición.} $F:I\times\mathbb{R}\times\mathbb{R}\to \mathbb{R}$ integrable, e indicaremos sus variables por $F = F(x,y,p)$. 

\begin{notacion}
    Usaremos los convenios:
    \begin{equation*}
        F_x = \frac{\partial F}{\partial x}, \qquad F_y = \frac{\partial F}{\partial y}, \qquad F_p = \frac{\partial F}{\partial p}
    \end{equation*}
    Y los análogos para derivadas parciales segundas:
    \begin{equation*}
        F_{yp} = \frac{\partial^2 F}{\partial y\partial p}
    \end{equation*}
\end{notacion}

\begin{lema}[Fundamental del cálculo de variaciones]
    Sea $I = [x_0,x_1]$, sea $f\in C^1(I)$. Si:
    \begin{equation*}
        \int_{x_0}^{x_1} f(x)\varphi(x)~dx = 0  \qquad \forall \varphi \in C_0(I)
    \end{equation*}
    Entonces $f(x) = 0 \quad \forall x\in I$.
    \begin{proof}
        Por contrarrecíproco, podemos suponer que $\exists \overline{x}\in \left]x_0,x_1\right[$ tal que $f(\overline{x})>0$. Así, existe un entorno de $\overline{x}$ tal que $f(x) > 0 \quad \forall x\in \left]\overline{x_0},\overline{x_1}\right[$. Podemos encontrar $\varphi \in C_0(I)$ tal que
        \begin{equation*}
            \int_{x_0}^{x_1} f(x)\varphi(x)~dx  > 0
        \end{equation*}
        Por ejemplo, podemos construir $\varphi$ a partir de la función:
        \begin{equation*}
            \rho(x) = \left\{\begin{array}{ll}
                e^{\frac{1}{|x|^2-1}} & \text{si\ } |x| < 1 \\
                 0 & \text{en otro caso\ } 
            \end{array}\right. 
        \end{equation*}
        donde $\rho \in C_0^\infty([-1,1])$.
    \end{proof}
\end{lema}

\begin{observacion}
    Varios comentarios:
    \begin{itemize}
        \item En el Lema fundamental se puede cambiar $C_0(I)$ por $C_0^1(I), \ldots, C_0^k(I), \ldots, C_0^\infty(I)$.
        \item Hay una versión alternativa del Lema, que a continuación se enuncia.
        \item De hecho, en dicho enunciado podemos también exigir para toda $\varphi \in C_c^\infty(I)$.
    \end{itemize}
\end{observacion}

\begin{lema}
    Sea $I = \left]x_0,x_1\right[$, $f\in L^1_{\text{loc}}(I)$. Si:
    \begin{equation*}
        \int_{x_0}^{x_1} f(x)\varphi(x)~dx = 0 \qquad \forall \varphi \in C_c(I)
    \end{equation*}
    Entonces, $f(x) = 0$ para casi todo $x\in I$.
    \begin{proof}
        El dual de $C_c(I)$ contiene a $L^1_{\text{loc}}(I)$ y se extrae a partir de ahí.
    \end{proof}
\end{lema}

\begin{teo}
    Sea $I=[x_0,x_1]$, $F\in C^2(I\times\mathbb{R}\times\mathbb{R})$. Consideramos el funcional
    \begin{equation*}
        \cc{F}(y) = \int_{x_0}^{x_1} F(x,y(x),y'(x))~dx 
    \end{equation*}

    definido sobre:
    \begin{equation*}
        D = \{y\in C^1([x_0,x_1]) : y(x_0) = y_0, y(x_1) = y_1\}, \quad y_0,y_1\in \mathbb{R}
    \end{equation*}
    Dado $y_\ast\in D\cap C^2(I)$ punto crítico de $\cc{F}$, satisface la ecuación de Euler-Lagrange:
    \begin{gather*}
        F_y(x,y_\ast(x),y_\ast'(x)) - \frac{d}{dx}\left[F_p(x,y_\ast(x),y_\ast'(x))\right] = 0 \qquad \forall x\in I \\
        \\
        \boxed{F_y - \frac{d}{dx}(F_p) = 0}
    \end{gather*}
    \begin{proof}
        Dado $y\in D$, $\varphi \in X\setminus \{0\}$, notaremos $y_\varepsilon = y+\varepsilon\varphi$, calculemos $\delta_\varphi\cc{F}(y)$ formalmente:
        \begin{gather*}
            \frac{d y_\varepsilon}{d \varepsilon} = \varphi \\
            y_\varepsilon' = y' + \varepsilon\varphi' \Longrightarrow \frac{d y'_\varepsilon}{d \varepsilon} = \varphi'
        \end{gather*}
        De donde:
        \begin{align*}
            \frac{d }{d \varepsilon} \cc{F}(y_\varepsilon) &= \frac{d }{d \varepsilon}\int_{x_0}^{x_1} F(x,y_\varepsilon(x),y_\varepsilon'(x))~dx  \\ 
                                                           &= \int_{x_0}^{x_1} F_y(x,y_\varepsilon(x),y_\varepsilon'(x))\frac{d y_\varepsilon}{d \varepsilon} + F_p(x,y_\varepsilon(x),y_\varepsilon'(x))\frac{dy_\varepsilon'}{d\varepsilon}~dx  \\
                                                           &= \int_{x_0}^{x_1} F_y(x,y_\varepsilon(x),y_\varepsilon'(x))\varphi(x) + F_p(x,y_\varepsilon(x),y_\varepsilon'(x))\varphi'(x)~dx 
        \end{align*}
        Por lo que:
        \begin{equation*}
            \left[\frac{d}{d\varepsilon}\cc{F}(y_\varepsilon)\right]_{\big|\varepsilon=0} = \int_{x_0}^{x_1} F_y(x,y(x),y'(x))\varphi(x) + F_p(x,y(x),y'(x))\varphi'(x)~dx 
        \end{equation*}
        Así:
        \begin{equation*}
            \delta_\varphi \cc{F}(y) = \int_{x_0}^{x_1} F_y\varphi + F_p\varphi'~dx 
        \end{equation*}
        Nos queremos quitar $\varphi'$ de esta expresión, lo conseguiremos integrando por partes:
        \begin{multline*}
            \int_{x_0}^{x_1} F_p(x,y(x),y'(x))\varphi'(x)~dx  = -\int_{x_0}^{x_1} \frac{d}{dx}[F_p(x,y(x),y'(x))]\varphi(x)~dx \\+ F_p(x_1,y(x_1),y'(x_1))\varphi(x_1) - F_p(x_0,y(x_0),y'(x_0))\varphi(x_0)
        \end{multline*}
        De donde:
        \begin{equation*}
            \delta_\varphi \cc{F}(y) = \int_{x_0}^{x_1} \left(F_y - \frac{d}{dx}F_p\right)\varphi~dx  + [F_p\varphi]_{x_0}^{x_1}
        \end{equation*}
        ¿Cuáles son las direcciones admisibles? Dada $y\in D$, ¿cómo ha de ser $\varphi$ para que $y_\varepsilon = y+\varepsilon\varphi \in D$? Tomando $\varphi \in C^1([x_0,x_1])$ tenemos la condición de regularidad, y vemos que:
        \begin{align*}
            y_\varepsilon(x_0) &= y(x_0) + \varepsilon\varphi(x_0) = y_0 \quad\Longleftrightarrow\quad \varphi(x_0) = 0
        \end{align*}
        Análogamente necesitamos que $\varphi(x_1) = 0$. Por tanto, la direcciones admisibles son $\varphi \in C^1_0([x_0,x_1])$. Usando estas direcciones admisibles, vemos que:
        \begin{equation*}
            \delta_\varphi \cc{F}(y) = \int_{x_0}^{x_1} \left(F_y - \frac{d}{dx}F_p\right)\varphi~dx  + \cancelto{0}{[F_p\varphi]_{x_0}^{x_1}}
        \end{equation*}
        Para cualquier dirección admisible $\varphi$. Usando que $y_\ast$ es punto crítico, tenemos:
        \begin{equation*}
            \delta_\varphi \cc{F}(y_\ast) = \int_{x_0}^{x_1}\left( F_y(x,y_\ast(x),y_\ast'(x)) - \frac{d}{dx}[F_p(x,y_\ast(x),y_\ast'(x))]\right)\varphi(x)~dx  = 0
        \end{equation*}
        Para toda $\varphi \in C_0^1(I)$. Por tanto, por el Lema fundamental del cálculo de variaciones tenemos la ecuación de Euler-Lagrange.
    \end{proof}
\end{teo}

\noindent
Si los candidatos no son de clase 2, podremos definir soluciones de la ecuación de Euler-Lagrange a partir de la fórmula:
\begin{equation*}
    \delta_\varphi\cc{F}(y) = \int_{x_0}^{x_1} F_y\varphi + F_p \varphi'~dx 
\end{equation*}
Pero involucra conceptos de análisis funcional más finos.

\begin{ejemplo}
    Estudiar las curvas planas que unen dos puntos dados con mínima longitud.\\

    \noindent
    Tomamos $A = (x_0,y_0), B = (x_1,y_1) \in \mathbb{R}^2$ y supondremos que $x_0<x_1$. Vamos a considerar únicamente aquellas curvas que pueden ser descritas como gráficas de una función\footnote{Pues nos parece que las curvas que optimizan el problema son de este tipo y simplifica mucho el problema.}, es decir, aplicaciones $x\longmapsto (x,y(x))$ con:
    \begin{equation*}
        y \in  D = \{f\in C^1([x_0,x_1]) : f(x_0) = y_0, f(x_1) = y_1\}
    \end{equation*}
    Así, las curvas tienen extremos en los puntos $A$ y $B$. Para este tipo de curvas se calcula su longitud como:
    \begin{equation*}
        L(y) = \int_{x_0}^{x_1} \|(1,y'(x))\|_2~dx = \int_{x_0}^{x_1} \sqrt{1+{(y'(x))}^{2}}~dx 
    \end{equation*}
    Así, tenemos $F= F(p) = \sqrt{1+p^2}$. La ecuación de Euler-Lagrange en este caso es:
    \begin{gather*}
        F_p = \frac{p}{\sqrt{1+p^2}} \\
        0 = F_y - \frac{d}{dx}(F_p) = -\frac{d}{dx} \left(\frac{y'(x)}{\sqrt{1+{(y'(x))}^{2}}}\right) \qquad \forall x\in [x_0,x_1]
    \end{gather*}
    Por lo que todos los puntos críticos $y(x)$ de $L$ satisfacen esta ecuación diferencial, de donde podemos deducir que para cierto $c_1\in \mathbb{R}$:
    \begin{equation*}
        \frac{y'(x)}{\sqrt{1+{(y'(x))}^{2}}} = c_1  \quad\Longrightarrow\quad y'(x) = c_1\sqrt{1+{(y'(x))}^{2}}
    \end{equation*}
    
    de donde:
    \begin{equation*}
        {(y'(x))}^{2} = c_1 \left( 1+{(y'(x))}^{2} \right) \quad\Longleftrightarrow\quad {(y'(x))}^{2} = c_2
    \end{equation*}
    para cierto $c_2\in \mathbb{R}$. Como $y\in C^1([x_0,x_1])$, tenemos entonces que $y'(x)$ es constante $\forall x\in [x_0,x_1]$, con lo que $y(x) = c_3 + c_4x\quad \forall x\in [x_0,x_1]$.\\

    \noindent
    Las constantes $c_3$ y $c_4$ se determinan por las condiciones $y(x_i) = y_i$, $i \in \{0,1\}$, obteniendo un único \underline{candidato a solución}, que se corresponde con la línea recta. Habría que comprobar que este candidato es efectivamente un mínimo de la función a optimizar.
\end{ejemplo}

\begin{definicion}
    Diremos que $y\in X$ es un \textbf{extremal} del funcional $\cc{F}$ si resuelve la ecuación de Euler-Lagrange asociada (notemos que no se pide $y\in D$).
\end{definicion}

\section{Problemas con frontera libre}
\noindent
En condiciones como las del anterior resultado, supongamos que 
\begin{equation*}
    D = \{y\in C^1([x_0,x_1]) : y(x_0) = y_0\}, \qquad y_0\in \mathbb{R}
\end{equation*}
Todo punto crítico $y_\ast \in D\cap C^2(I)$, además de cumplir la ecuación de Euler-Lagrange, cumple la condición extra\footnote{La estructura del problema lo exige, pues quitamos una condición y esta tiene que estar por otra parte para tener solución a una EDO de orden 2.} $F_p(x_1,y_\ast(x_1), y_\ast'(x_1)) = 0$.\\

\noindent
Esto se debe a que:
\begin{equation*}
    \delta_\varphi \cc{F}(y) = \int_I \left(F_y - \frac{d}{dx}(F_p)\right)\varphi~dx + \left[F_p(x,y(x),y'(x))\varphi(x)\right]_{x_0}^{x_1}
\end{equation*}
Si $\varphi \in C_0^1(I)$ son admisibles, entonces Euler-Lagrange se cumple, y el término asociado a $x=x_0$ se va a cero por la condición del dominio $D$.
